% Source: https://artofproblemsolving.com/community/c3h1925244_calculate_the_area_of_a_triangle
% Stable identity: aops-topic-1925244-q01
Let $T$ be the area of triangle $ABC$. Since $F$ is the midpoint of $AC$, $[BFC]=\dfrac {T}{2}$.

\begin{asy}
usepackage("amsmath");
size(6.5cm);
defaultpen(fontsize(10pt));
pair B=(0,0),C=(6,0),A=(1.8,4.2),D=(1,0),Ee=(4,0);pair F=(A+C)/2;
pair G=extension(B,F,A,D),H=extension(B,F,A,Ee);
fill(D--G--H--Ee--cycle,lightgray);
draw(A--B--C--cycle,linewidth(1));draw(A--D);draw(A--Ee);draw(B--F);
dot(D);dot(Ee);dot(F);dot(G);dot(H);
label("$A$",A,N);label("$B$",B,SW);label("$C$",C,SE);label("$D$",D,S);label("$E$",Ee,S);
label("$F$",F,NE);label("$G$",G,NW);label("$H$",H,N);
label("$51$",(G+H+D+Ee)/4+(0.15,-0.1));
label("$k$",(B+D)/2,S);label("$3k$",(D+Ee)/2,S);label("$2k$",(Ee+C)/2,S);
label("\mbox{}",(min(currentpicture,user=true)+max(currentpicture,user=true))/2,black); // keeps the text after the figure black when LaTeX typesets the labels inline
\end{asy}

\textbf{Locating $G$ and $H$ on $BF$ (mass points).} For $G=AD\cap BF$: give $A$ and $C$ mass $1$ each, so $F$ (the midpoint) has mass $2$. Since $BD:DC=1:5$, $B$ needs mass $5$, so that $D$ balances $B$ and $C$. Then $G$ balances $B$ (mass $5$) and $F$ (mass $2$): $BG:GF=2:5$, so $BG=\dfrac{2}{7}BF$.

For $H=AE\cap BF$: now $BE:EC=4:2=2:1$. Give $A$ and $C$ mass $2$, so $F$ has mass $4$, and $B$ mass $1$. Then $BH:HF=4:1$, so $BH=\dfrac{4}{5}BF$.

\textbf{Areas.} Triangles with a vertex at $B$ and sides along $BC$ and $BF$ scale with both sides:
\[
[BHE]=\dfrac{BE}{BC}\cdot\dfrac{BH}{BF}\cdot[BCF]=\dfrac{4}{6}\cdot\dfrac{4}{5}\cdot\dfrac {T}{2}=\dfrac{4T}{15},\qquad
[BGD]=\dfrac{BD}{BC}\cdot\dfrac{BG}{BF}\cdot[BCF]=\dfrac{1}{6}\cdot\dfrac{2}{7}\cdot\dfrac {T}{2}=\dfrac{T}{42}.
\]
The quadrilateral is the difference:
\[
[DGHE]=[BHE]-[BGD]=\dfrac{4T}{15}-\dfrac{T}{42}=\dfrac{56T-5T}{210}=\dfrac{17T}{70}.
\]
So $\dfrac{17T}{70}=51$ and $T=\boxed{210}$.
